\documentclass[a4paper]{article}

% --- PACHETE ESENȚIALE PENTRU COMPACTARE ---
\usepackage[a4paper, margin=0.2in]{geometry} % set tight margins
\usepackage{multicol}      % create a multi-column layout
\usepackage{titlesec}      % customize section/subsection titles
\usepackage{paralist}      % creating compact lists
\usepackage{amsmath}       % math environments
\usepackage[romanian]{babel} % Limba romana
\usepackage[utf8]{inputenc}  % Encoding
\usepackage[T1]{fontenc}     % Font encoding

% --- PACHETE UTILITARE & STIL ---
\usepackage{xcolor}        % custom colors
\usepackage{graphicx}      % images
\usepackage{listings}      % code snippets
\usepackage{hyperref}      % clickable links
\usepackage{tcolorbox}     % Pentru definitii (daca e nevoie)
\usepackage{booktabs}      % Adaugat pentru tabele frumoase
\usepackage{array}         % Pentru coloane p{width} in tabele

% --- DEFINIȚII CULORI ---
\definecolor{codegray}{rgb}{0.9,0.9,0.9} % code background
\definecolor{darkblue}{rgb}{0.0,0.0,0.5} % for comments
\definecolor{defbox_bg}{rgb}{0.95, 0.95, 1.0} % definition box background
\definecolor{defbox_fr}{rgb}{0.4, 0.4, 0.8} % frame

% --- CONFIGURARE GLOBALĂ LAYOUT & SPAȚIERE ---

% set paragraph spacing: no indent, with tiny vertical gap between paragraphs
\setlength{\parindent}{0pt}
\setlength{\parskip}{0.5ex}

% make section and subsection titles extremely compact
\titlespacing*{\section}{0pt}{0.95ex}{0.7ex}
\titlespacing*{\subsection}{0pt}{0.8ex}{0.4ex}
\titleformat*{\section}{\bfseries\small}
\titleformat*{\subsection}{\bfseries\footnotesize}

% make lists as compact as possible
\let\itemize\compactitem
\let\enditemize\endcompactitem
\let\enumerate\compactenum
\let\endenumerate\endcompactenum

% ****** SOLUȚIA PENTRU SPAȚII ******
\raggedbottom % Nu mai echilibra vertical coloanele/paginile

% --- STIL BLOCURI COD ---
\lstdefinestyle{mystyle}{
    backgroundcolor=\color{codegray},   
    commentstyle=\color{darkblue},
    keywordstyle=\color{blue},
    numberstyle=\tiny\color{gray},
    stringstyle=\color{red},
    basicstyle=\ttfamily\tiny,
    breaklines=true,
    numbers=left,                    
    numbersep=4pt,                  
    showspaces=false,                
    showstringspaces=false,
    showtabs=false,                  
    tabsize=2,
    mathescape=false % Asigura ca $ e tratat ca text
}
\lstset{style=mystyle}

% --- LIMBAJ x86 ---
\lstdefinelanguage{x86asm}{
  keywords={mov, movl, add, addl, sub, subl, imul, mul, div, idiv, idivl, cdq, int, jmp, je, jne, jg, jl, jge, jle, ja, jb, jae, jbe, cmp, cmpl, lea, leal, loop, pushl, popl, call, ret, .global, .text, .data, .asciz, .ascii, .long, .word, .byte, .space, flds, fstps, fldl, fstpl, movss, addss, subss, mulss, divss, sqrtss, cvtsi2ss, cvtss2sd, maxss, minss},
  morecomment=[l]{\#}, 
  sensitive=false
}

% --- LIMBAJ RISC-V ---
\lstdefinelanguage{RISCV}{
  keywords={lw, sw, lb, sb, lh, sh, add, sub, mul, div, rem, and, or, xor, sll, srl, addi, slti, andi, lui, auipc, li, la, mv, nop, jal, j, jr, ret, beq, bne, blt, bge, bltu, bgeu, beqz, bnez, ecall, sw, lw, addi, c.lwsp, c.swsp, c.addi, c.j, c.beqz, c.bnez},
  morecomment=[l]{\#},
  sensitive=false
}

% --- HYPERREF SETUP---
\hypersetup{
    colorlinks=true, linkcolor=blue, urlcolor=cyan,
}

% --- COMANDĂ CUSTOM PENTRU DEFINIȚII ---
% Simpler alternative w/o tcolorbox (din template-ul tau)
\newcommand{\definition}[2]{\par\noindent\fcolorbox{defbox_fr}{defbox_bg}{\begin{minipage}{0.95\columnwidth}\textbf{#1:} #2\end{minipage}}\par}

\begin{document}

% --- CUSTOM, MINIMAL TITLE ---
\noindent
\textbf{\tiny Cheat Sheet - Curs ASC (x86 vs RISC-V)}
\hrule
\vspace{0.2ex}

% --- GLOBAL SETTINGS FOR THE DOCUMENT BODY ---
\pagestyle{empty}
\footnotesize

\begin{multicols}{2} % Începe layout-ul pe două coloane

\section{x86 AT\&T (Arhitectura CISC)}

\subsection{Registri și Notație}

\definition{Registri GPR (32-bit)}{
    \textbf{EAX} (Accumulator), \textbf{EBX} (Base), \textbf{ECX} (Counter), \textbf{EDX} (Data), \textbf{ESP} (Stack Pointer), \textbf{EBP} (Base Pointer), \textbf{ESI} (Source Index), \textbf{EDI} (Destination Index).
    Parțile inferioare: \textbf{AX} (16-bit), \textbf{AH} (high 8-bit), \textbf{AL} (low 8-bit).
}

\definition{EFLAGS (Indicatori)}{
    \textbf{ZF} (Zero Flag): Setat dacă rezultatul e 0.
    \textbf{SF} (Sign Flag): Setat dacă rezultatul e negativ.
    \textbf{CF} (Carry Flag): Setat dacă a avut loc transport.
    \textbf{OF} (Overflow Flag): Setat la depășire cu semn.
}

\definition{Sintaxa AT\&T}{
    Sursa \textit{înaintea} Destinației: \texttt{movl \%eax, \%ebx}
}

\definition{Prefixe Sintaxă}{
    Registri: \texttt{\%eax}. Valori imediate: \texttt{\$10}. Etichete/Adrese: \texttt{et}/\texttt{\$et}.
}

\definition{Adresare Memorie}{
    \texttt{offset(\%baza, \%index, scalar)} $\rightarrow$ \texttt{\%baza + \%index * scalar + offset}
}

\definition{Tipuri Date}{
    \texttt{.byte} (1B), \texttt{.word} (2B), \texttt{.long} (4B), \texttt{.ascii} (fără null), \texttt{.asciz} (cu null), \texttt{.space N} (N bytes).
}

\subsection{Apeluri de Sistem (Linux 32-bit)}

\definition{Apel (Interrupt)}{
    Se folosește \texttt{int \$0x80}. Codul funcției $\rightarrow$ \texttt{\%eax}. Argumente $\rightarrow$ \texttt{\%ebx, \%ecx, \%edx}.
}

\begin{lstlisting}[language=x86asm]
# sys_exit(0) - Opreste programul
movl $1, %eax
movl $0, %ebx
int $0x80
\end{lstlisting}
\begin{lstlisting}[language=x86asm]
# sys_write(stdout, msg, len) - Afiseaza
movl $4, %eax
movl $1, %ebx
movl $msg_addr, %ecx
movl $msg_len, %edx
int $0x80
\end{lstlisting}

% --- SECȚIUNEA EXTINSĂ ÎNCEPE AICI ---

\subsection{Stiva (Stack) și Proceduri}

\definition{Stiva (LIFO)}{
    Crește spre adrese mici. \texttt{\%esp} (Stack Pointer) indică vârful stivei. \texttt{\%ebp} (Base Pointer) indică baza cadrului de stivă curent (pentru referință stabilă).
}

\definition{Instrucțiuni Stivă}{
    \begin{compactitem}
        \item \texttt{pushl src} (Doar \texttt{long} (4B) sau \texttt{word} (2B)). Decrementează \texttt{\%esp} cu 4, apoi scrie \texttt{src} la \texttt{(\%esp)}.
        \item \texttt{popl dest} (Doar \texttt{long} sau \texttt{word}). Citește valoarea de la \texttt{(\%esp)} în \texttt{dest}, apoi incrementează \texttt{\%esp} cu 4.
    \end{compactitem}
}

\definition{Apel Procedură}{
    \begin{compactitem}
        \item \texttt{call etic}: 1. Pune adresa instrucțiunii următoare (adresa de retur) pe stivă. 2. Sare la \texttt{etic}.
        \item \texttt{ret}: 1. Ia adresa de retur de pe stivă. 2. Sare la acea adresă.
    \end{compactitem}
}

\definition{Convenție Apel C (Argumente)}{
    \begin{compactitem}
        \item Argumentele sunt puse pe stivă de \textit{apelant}, în ordine \textbf{inversă}.
        \item \textit{Apelantul} curăță stiva după apel (ex: \texttt{addl \$8, \%esp} pentru 2 argumente \texttt{long}).
    \end{compactitem}
}

\begin{lstlisting}[language=x86asm]
# Prolog Procedura (Setup Stack Frame)
eticheta_proc:
    pushl %ebp       # 1. Salveaza vechiul %ebp
    movl %esp, %ebp  # 2. Seteaza noul %ebp la %esp curent
    subl $8, %esp    # 3. Aloca spatiu pt 2 var locale
\end{lstlisting}

\begin{lstlisting}[language=x86asm]
# Epilog Procedura (Restore Stack Frame)
    movl %ebp, %esp  # 1. Elibereaza var locale (sau leave)
    popl %ebp        # 2. Restaureaza vechiul %ebp
    ret              # 3. Revine la apelant
\end{lstlisting}

\definition{Layout Cadru Stivă (după prolog)}{
    Accesul se face relativ la \texttt{\%ebp}:
    \begin{compactitem}
        \item \texttt{...}
        \item \texttt{12(\%ebp)}: Al doilea argument
        \item \texttt{8(\%ebp)}: Primul argument
        \item \texttt{4(\%ebp)}: Adresa de retur (salvată de \texttt{call})
        \item \texttt{0(\%ebp)}: Vechiul \texttt{\%ebp} (salvat de prolog) $\leftarrow$ \textbf{\%ebp}
        \item \texttt{-4(\%ebp)}: Prima variabilă locală
        \item \texttt{-8(\%ebp)}: A doua variabilă locală
        \item \texttt{...} $\leftarrow$ \textbf{\%esp}
    \end{compactitem}
}

\definition{Salvare Registri}{
    \textbf{Caller-saved} (apelantul salvează dacă are nevoie de ele): \texttt{\%eax, \%ecx, \%edx}.
    \textbf{Callee-saved} (apelatul salvează dacă le folosește): \texttt{\%ebx, \%esi, \%edi}.
}

\definition{Returnare Valori}{
    Valorile sunt returnate de procedură (callee) prin registri, în această ordine: \textbf{\%eax} (prima valoare), \textbf{\%ecx} (a doua), \textbf{\%edx} (a treia).
}

% --- SECȚIUNEA EXTINSĂ SE TERMINĂ AICI ---

\subsection{FPU și SSE}

\definition{FPU (Stack-based)}{
    Folosește stiva \texttt{st(0)..st(7)}.
    \texttt{flds mem} (load float).
    \texttt{fstps mem} (store float + pop).
}

\definition{SSE (Register-based)}{
    Folosește registri \texttt{\%xmm0..\%xmm7}.
    \texttt{movss s, d} (mută float).
    \texttt{addss s, d} (adună float).
    \texttt{cvtsi2ss i, f} (int $\rightarrow$ float).
    \texttt{cvtss2sd f, d} (float $\rightarrow$ double).
}

\definition{printf float}{
    Argumentele float trebuie promovate la double (8B). Se folosește \texttt{cvtss2sd \%xmm0, \%xmm0} înainte de apel.
}

\subsection{Instrucțiuni Aritmetice / Logice}

\definition{Aritmetice}{
    \texttt{addl s, d} (d=d+s), \texttt{subl s, d} (d=d-s).
    \texttt{imul src} (EDX:EAX = EAX * src).
    \texttt{idivl src} (EAX = cat, EDX = rest).
}

\definition{Logice / Shift}{
    \texttt{andl s, d}, \texttt{orl s, d}, \texttt{xorl s, d}, \texttt{notl d}.
    \texttt{shll N, d}, \texttt{shrl N, d} (logic).
    \texttt{sall N, d}, \texttt{sarl N, d} (aritmetic).
}

\subsection{Salturi Condiționate x86 (după \texttt{cmp op1, op2})}
\tiny
\begin{tabular}{ll}
\toprule
\textbf{Cu Semn} & \textbf{Descriere (op2 ? op1)} \\
\midrule
\texttt{je} & jump if equal (==) \\
\texttt{jne} & jump if not equal (!=) \\
\texttt{jg} & jump if greater (>) \\
\texttt{jl} & jump if less (<) \\
\texttt{jge} & jump if greater or equal (>=) \\
\texttt{jle} & jump if less or equal (<=) \\
\midrule
\textbf{Fără Semn} & \textbf{Descriere (op2 ? op1)} \\
\midrule
\texttt{ja} & jump if above (>) \\
\texttt{jb} & jump if below (<) \\
\texttt{jae} & jump if above or equal (>=) \\
\texttt{jbe} & jump if below or equal (<=) \\
\bottomrule
\end{tabular}


\columnbreak % Trecere la a doua coloană

\section{RISC-V (Arhitectura RISC)}

\definition{Filosofie}{
    Arhitectură \textbf{open-source}, \textbf{modulară} (bază + extensii).
}

\subsection{Registri (RV32I)}

\definition{32 Registri GPR (x0-x31)}{
    \textbf{Nume ABI (Convenție):}
    \begin{compactitem}
        \item \texttt{x0 (zero)}: Hardcodat la 0.
        \item \texttt{x1 (ra)}: Adresa de retur.
        \item \texttt{x2 (sp)}: Stack Pointer.
        \item \texttt{x8 (s0/fp)}: Frame Pointer.
        \item \texttt{a0-a7 (x10-x17)}: Argumente / Valori retur.
        \item \texttt{s0-s11 (x8-9, x18-27)}: Callee-saved.
        \item \texttt{t0-t6 (x5-7, x28-31)}: Caller-saved.
    \end{compactitem}
}

\subsection{Arhitectura Load-Store}

\definition{Principiu}{
    Operațiile (add, sub etc.) se fac \textbf{doar} între registri. Memoria se accesează \textbf{doar} cu instrucțiuni \texttt{load} și \texttt{store}.
}

\definition{Acces Memorie}{
    \texttt{lw rd, offset(rs1)}: Load Word (4B).
    \texttt{sw rs2, offset(rs1)}: Store Word (4B).
    \texttt{lb/sb} (1B), \texttt{lh/sh} (2B).
}

\definition{Pseudo-Instrucțiuni}{
    \texttt{li rd, imm} (Load Immediate).
    \texttt{la rd, eticheta} (Load Address).
    \texttt{mv rd, rs} (Move, echivalent \texttt{addi rd, rs, 0}).
    \texttt{nop} (No operation).
}

\subsection{Apeluri de Sistem (ecall)}

\definition{Apel (ecall)}{
    Se folosește instrucțiunea \texttt{ecall}. Codul funcției $\rightarrow$ \texttt{a7} (\texttt{x17}). Argumente $\rightarrow$ \texttt{a0-a6} (\texttt{x10-x16}).
}

\begin{lstlisting}[language=RISCV]
# exit(0) - Opreste programul (Ripes)
li a7, 93
li a0, 0
ecall
\end{lstlisting}

\begin{lstlisting}[language=RISCV]
# PrintString(a0=adresa) (Ripes)
la a0, eticheta_string
li a7, 4
ecall
\end{lstlisting}

\subsection{Salturi și Proceduri}

\definition{Salturi}{
    \texttt{j eticheta} (Salt necondiționat, alias pt \texttt{jal x0, et}).
    \texttt{beq r1, r2, et} (Branch if equal).
    \texttt{bne, blt, bge} (Not equal, less than, greater/equal).
}

\definition{Apel Procedură}{
    \texttt{call eticheta} (Alias pt \texttt{jal ra, eticheta}). Salvează PC+4 în \texttt{ra} (\texttt{x1}).
    \texttt{ret} (Alias pt \texttt{jalr x0, ra, 0}).
}

\definition{Argumente (Convenție)}{
    Primele 8 argumente $\rightarrow$ \texttt{a0-a7}. Restul $\rightarrow$ pe stivă. Valori retur $\rightarrow$ \texttt{a0, a1}.
}

\begin{lstlisting}[language=RISCV]
# Prolog Procedura
addi sp, sp, -8   # aloca spatiu
sw ra, 4(sp)      # salveaza adresa retur
sw s0, 0(sp)      # salveaza frame pointer
addi s0, sp, 8    # seteaza noul fp
\end{lstlisting}

\begin{lstlisting}[language=RISCV]
# Epilog Procedura
lw s0, 0(sp)      # restaureaza fp
lw ra, 4(sp)      # restaureaza adr retur
addi sp, sp, 8    # elibereaza spatiu
ret
\end{lstlisting}

\subsection{Instrucțiuni RISC-V (RV32I/M)}
\tiny
\begin{tabular}{p{1.2cm} p{1.8cm} p{4.6cm}}
\toprule
\textbf{Instr.} & \textbf{Operanzi} & \textbf{Descriere} \\
\midrule
\texttt{add/sub} & rd, rs1, rs2 & rd = rs1 +/- rs2 \\
\texttt{mul/div} & rd, rs1, rs2 & Inmultire/Impartire (M) \\
\texttt{rem} & rd, rs1, rs2 & Rest (M) \\
\texttt{and/or/xor} & rd, rs1, rs2 & Logice pe biti \\
\texttt{addi} & rd, rs1, imm & rd = rs1 + imm \\
\texttt{slti} & rd, rs1, imm & rd = (rs1 < imm) ? 1 : 0 \\
\texttt{lui} & rd, imm & rd = imm << 12 \\
\texttt{auipc} & rd, imm & rd = pc + (imm << 12) \\
\bottomrule
\end{tabular}

\subsection{Salturi Condiționate RISC-V}
\tiny
\begin{tabular}{p{1.2cm} p{1.8cm} p{4.6cm}}
\toprule
\textbf{Instr.} & \textbf{Operanzi} & \textbf{Descriere (Salt dacă...)} \\
\midrule
\texttt{beq} & rs1, rs2, imm & rs1 == rs2 \\
\texttt{bne} & rs1, rs2, imm & rs1 != rs2 \\
\texttt{blt} & rs1, rs2, imm & rs1 < rs2 (cu semn) \\
\texttt{bge} & rs1, rs2, imm & rs1 >= rs2 (cu semn) \\
\texttt{bltu} & rs1, rs2, imm & rs1 < rs2 (fără semn) \\
\texttt{bgeu} & rs1, rs2, imm & rs1 >= rs2 (fără semn) \\
\bottomrule
\end{tabular}

\subsection{Instrucțiuni RVC (Comprimate)}
\tiny
\begin{tabular}{p{1.2cm} p{2.1cm} p{4.3cm}}
\toprule
\textbf{Instr.} & \textbf{Operanzi} & \textbf{Descriere} \\
\midrule
\texttt{c.addi} & rd, imm & rd = rd + imm (imm != 0) \\
\texttt{c.lwsp} & rd, offset(sp) & Load Word (din stiva) \\
\texttt{c.swsp} & rd, offset(sp) & Store Word (pe stiva) \\
\texttt{c.j} & offset & Salt neconditionat \\
\texttt{c.beqz} & rs, offset & Salt daca rs == 0 \\
\texttt{c.bnez} & rs, offset & Salt daca rs != 0 \\
\texttt{c.lw/c.sw} & rd', off(rs') & Load/Store (pt registri s0-s1, a0-a5) \\
\bottomrule
\end{tabular}

\section{Arhitectură (Pipeline \& Cache)}
\definition{Pipeline 5 Etape}{
    \textbf{F}etch (Ia instrucțiune).
    \textbf{D}ecode (Decodează, citește registri).
    \textbf{E}xecute (Calculează în ALU).
    \textbf{M}emory (Accesează memoria).
    \textbf{W}rite-Back (Scrie rezultatul în registru).
}

\definition{Hazard-uri}{
    \textbf{De Date}: O instrucțiune are nevoie de rezultatul uneia nefinalizate.
    \textbf{De Control}: Salt (branch) - nu se știe următoarea instrucțiune.
    \textbf{Structurale}: Două instrucțiuni cer aceeași resursă hardware.
}

\definition{Soluții Hazard}{
    \textbf{Stalling}: Introducere \texttt{nop} (bule) în pipeline.
    \textbf{Forwarding}: Rezultat ALU trimis direct la intrarea ALU, ocolind scrierea în registru.
    \textbf{Branch Prediction}: Ghicirea direcției saltului.
}

\definition{Localitate (Cache)}{
    \textbf{Temporală}: Date accesate recent vor fi accesate iar.
    \textbf{Spațială}: Date apropiate (ex. în array) vor fi accesate curând.
}

\definition{Mapare Cache}{
    \textbf{Direct-Mapped}: Adresă $\rightarrow$ o singură linie specifică (index).
    \textbf{N-Way Set Associative}: Adresă $\rightarrow$ oricare din N linii dintr-un set.
    \textbf{Fully Associative}: Adresă $\rightarrow$ orice linie din cache.
}

\end{multicols}

\end{document}